\section{Introduction}

Let $A$ be a small category, non-empty and connected. Four properties of $A$
are compared in this paper; the definitions are in
Section~\ref{sec:classes}.
\begin{itemize}
\item $A$ is \emph{filtered}.
\item $A$ is \emph{exact}: the colimit functor from $A$-indexed diagrams of
 abelian groups to abelian groups is exact.
\item $A$ is \emph{total}: every category that receives a final functor from
 $A$ is weakly contractible. This is Grothendieck's \emph{asph\'ericit\'e
 totale} \cite{Der}; it holds if and only if every connected component of the
 category of elements of every functor $A\to\Set$ is weakly contractible
 (Proposition~\ref{prop:total}), and if and only if colimits over $A$ of
 diagrams of simplicial sets are homotopy colimits (Remark~\ref{rem:hocolim}).
\item $A$ is \emph{quasi-filtered}: pairs of arrows with a common source can be
 completed to commutative squares, and parallel pairs have affine equalizers.
\end{itemize}
Each property implies the next (Propositions~\ref{prop:incl}
and~\ref{prop:qf}). Isbell and Mitchell showed in 1973 that the first
implication cannot be reversed. In the manuscript \emph{Les D\'erivateurs}
Grothendieck asks whether total categories are filtered. The results of this
paper are the following.

\begin{theorem}\label{thm:A}
Let $\Cd$ be the category of finite linearly ordered $3$-uniform hypergraphs
in which every non-empty set $U$ of vertices satisfies $2|U|-3e(U)>0$, where
$e(U)$ is the number of hyperedges contained in $U$, with the closed
embeddings as arrows \textup{(}Definition~\ref{def:cdelta}\textup{)}. Then $\Cd$ is
total, and for every non-zero ring $R$ the colimit functor on $\Cd$-indexed
diagrams of $R$-modules is not exact.
\end{theorem}

\begin{theorem}\label{thm:B}
The category $\mathcal T$ of finite ordered two-graphs and their embeddings
is quasi-filtered and is not total.
\end{theorem}

\begin{theorem}\label{thm:C}
Among non-empty connected small categories,
\[
 \{\text{filtered}\}\subsetneq\{\text{exact}\}\subsetneq\{\text{total}\}
 \subsetneq\{\text{quasi-filtered}\}.
\]
The three inclusions are strict by the category $\Delta_{\mathrm{inj}}$ of
the finite non-empty ordinals and increasing injections, by $\Cd$, and by
$\mathcal T$. Moreover a non-empty connected category is filtered if and
only if it is total and sifted.
\end{theorem}

Theorem~\ref{thm:A} is Theorems~\ref{thm:total} and~\ref{thm:notexact};
Theorem~\ref{thm:B} is Theorem~\ref{thm:twographs}. For
Theorem~\ref{thm:C}: the inclusions are Propositions~\ref{prop:incl}
and~\ref{prop:qf}; $\Delta_{\mathrm{inj}}$ is exact by
\cite[Theorems~1 and~3]{IM73}, and it is not filtered, because no arrow
equalizes the two arrows from the ordinal with one element to the ordinal
with two; the last assertion is Proposition~\ref{prop:sifted}.

Theorem~\ref{thm:A} separates two weakenings of filteredness that are both
classical: exactness of colimits of modules, and the property that colimits of
spaces are homotopy colimits. The proofs of total asphericity
(Theorem~\ref{thm:A}) and of its failure (Theorem~\ref{thm:B}) follow one
pattern. Both categories consist of the finite closed substructures of a
countable homogeneous structure $M$. If every open subgroup of $\Aut(M)$ is
the pointwise stabilizer of a finite closed set, then every component of
every category of elements is weakly contractible or weakly equivalent to
the category itself (Proposition~\ref{prop:reduction}). So only the
classifying space of the category matters: it is contractible for $\Cd$, and
simply connected with non-zero second homology for $\mathcal T$. Exactness is
decided by a count: in a hypergraph of $\Cd$ any $r$ hyperedges span more
than $3r/2$ vertices, so one of them has a vertex of its own.

Section~\ref{sec:manuscript} relates the results to the passage of
\emph{Les D\'erivateurs}. This paper was produced by AI models; who proved,
wrote and checked what is recorded in Section~\ref{sec:provenance}.
